\documentclass[11pt]{article}
\usepackage[a4paper,margin=1in]{geometry}
\usepackage{amsmath,amssymb,mathtools,bm}
\usepackage{enumitem,booktabs,array}
\usepackage{tcolorbox}
\usepackage{hyperref}
\usepackage[protrusion=true,expansion=false]{microtype}
\hypersetup{colorlinks=true,linkcolor=blue,urlcolor=blue,citecolor=blue}
\setlist[itemize]{topsep=2pt,itemsep=2pt}
\setlist[enumerate]{topsep=2pt,itemsep=2pt}
\newtcolorbox{keybox}{colback=blue!3!white,colframe=blue!40!black,boxrule=0.4pt,arc=1mm}
\newtcolorbox{alertbox}{colback=red!3!white,colframe=red!40!black,boxrule=0.4pt,arc=1mm}
\newtcolorbox{answerbox}{colback=green!3!white,colframe=green!40!black,boxrule=0.4pt,arc=1mm}
\newtcolorbox{drillbox}{colback=yellow!5!white,colframe=orange!60!black,boxrule=0.4pt,arc=1mm}
\newcommand{\EV}{\mathbb{E}}
\newcommand{\blank}[1]{\underline{\hspace{#1}}}

\title{W07-01: G\'eczy, Minton, Schrand (1997, JF) \\
\large ``Why Firms Use Currency Derivatives'' --- Active Recall Workbook}
\author{Banking \& Risk Management --- Exam Prep}
\date{\today}
\begin{document}\maketitle

\begin{keybox}
\textbf{Paper in one sentence.} 372 Fortune~500 nonfinancial firms in 1990: 41\% use FX derivatives; logit with 75--78\% correct-classification shows usage rises with R\&D (underinvestment / growth options), falls with the quick ratio (internal liquidity), rises with size (scale economies in hedging), and \emph{foreign-denominated debt substitutes for FX derivatives as a natural hedge}.

\textbf{Placement.} Empirical sequel to Smith--Stulz (1985) and Froot--Scharfstein--Stein (1993). Direct cousin of Tufano (1996): same Tobit/logit style, but FX rather than gold, and focused on the \emph{decision} to hedge rather than the \emph{extent}.
\end{keybox}

\section*{Round 1: Complete Walk-Through (read carefully once)}

\subsection*{1.1 Sample and stylised facts}
\begin{itemize}
\item 372 Fortune~500 nonfinancial firms with \emph{ex-ante} FX exposure (foreign sales, foreign-denominated debt, foreign subsidiaries, or foreign income tax); fiscal year-end 1991 derivative disclosures.
\item 41.4\% use some FX derivative; 59.1\% use some derivative overall; usage is concentrated in Chemicals, Computers/Electronics, Pharmaceuticals.
\item Forwards are by far the most common instrument; swaps and options are secondary.
\end{itemize}

\subsection*{1.2 Hypotheses mapped to rationales}
\begin{tabular}{p{3.4cm}p{3.6cm}p{4.8cm}p{1.6cm}}
\toprule
\textbf{Variable} & \textbf{Proxy for} & \textbf{Rationale} & \textbf{Pred.} \\
\midrule
R\&D / sales & Growth options, \newline underinvestment (FSS) & Firms with high convex future investment want cash in bad states; hedge to avoid costly external finance & $+$ \\
Capex / sales & Investment opportunities & Same FSS logic & $+$ \\
Quick ratio & Internal liquidity & Already cash-rich firms can self-insure, less need to hedge & $-$ \\
Book leverage & Financial distress & Higher expected bankruptcy cost $\Rightarrow$ hedge & $+$ \\
Size (log assets) & Scale economies in hedging & Large fixed cost of running a hedging program & $+$ \\
Interest coverage & Distress (inverse) & Low coverage $\Rightarrow$ hedge more & $-$ \\
Foreign debt / assets & Natural hedge of foreign revenue & If you already have FX debt, you don't \emph{need} FX derivatives & $-$ \\
Foreign pretax income / sales & Exposure / natural hedge & Exposed firms hedge more, but only once you control for natural hedges & ambig. \\
Dividend payout & Internal funds & Mixed --- either distress proxy or commitment to pay out & $-/+$ \\
Tax convexity (NOL, ITC) & Smith--Stulz tax argument & Convexity of tax $\Rightarrow$ hedge to cut expected tax & $+$ \\
\bottomrule
\end{tabular}

\subsection*{1.3 Specification: logit for the hedge/no-hedge decision}
\[
\Pr(HEDGE_i=1)=\Lambda\bigl(\beta_0+\bm{\beta}^\top \mathbf{x}_i\bigr),\qquad \Lambda(z)=\frac{e^z}{1+e^z}.
\]
Marginal effect at the mean: $\partial\Pr/\partial x_k=\Lambda'(\bar z)\,\beta_k$, where $\Lambda'(\bar z)=\bar p(1-\bar p)$. GMS report the marginal effects $\Delta\!\Pr$ rather than raw $\beta_k$ (they are easier to interpret).

\subsection*{1.4 Headline results (Table~IV logit)}
\begin{itemize}
\item R\&D/sales: $\beta\approx 7$, marginal effect $\approx +2.5\%$ per 1 pp of R\&D --- \emph{highly} significant.
\item Quick ratio: $\beta\approx -1.9$, marginal effect $\approx -0.5$ --- significant negative.
\item Size (log assets): positive and significant --- supports scale economies.
\item Foreign debt: firms with foreign-denominated debt have materially \emph{lower} probability of using FX derivatives, controlling for exposure.
\item Tax variables (NOL carryforwards, ITCs): weak; they become meaningful mostly when interacted with income near the convex kink.
\item Correct-classification rate $\approx 75$--$78\%$.
\end{itemize}

\subsection*{1.5 Key identification / econometrics caveats}
\begin{alertbox}
\begin{itemize}
\item \textbf{Selection on exposure.} Sample conditions on having some FX exposure. Extending to all Fortune~500 would add mechanical zeros; GMS deliberately prune to get a cleaner decision problem.
\item \textbf{Endogeneity of leverage.} Leverage is a choice; simultaneous equations would be preferable (cf.\ Graham--Rogers W07-02).
\item \textbf{Intensity vs.\ extensive margin.} GMS estimate a binary hedge/no-hedge logit, so they can say \emph{who} hedges but not \emph{how much} --- contrast Tufano's fraction-of-production Tobit and Graham--Rogers' Tobit.
\item \textbf{Notional misclassification.} Pre-SFAS 119 disclosures are noisy; a firm reporting a forward of \$1 million and one reporting \$1 billion look identical on the binary DV.
\end{itemize}
\end{alertbox}

\section*{Round 2: Level 1 Fill-in (30\% missing)}
\begin{enumerate}
\item Sample: \blank{1.5cm} Fortune 500 nonfinancial firms with ex-ante FX exposure in fiscal year-end \blank{1.5cm}.
\item Among these, \blank{1cm}\% use FX derivatives.
\item The logit specifies $\Pr(HEDGE_i=1)=\Lambda(\beta_0+\bm\beta^\top \mathbf x_i)$ with $\Lambda(z)=\blank{2cm}$.
\item Marginal effect at the mean: $\partial\Pr/\partial x_k=\bar p(1-\bar p)\cdot\blank{1.5cm}$.
\item Sign predictions: R\&D/sales \blank{0.8cm}; quick ratio \blank{0.8cm}; size \blank{0.8cm}; foreign debt \blank{0.8cm}.
\item Correct-classification rate roughly \blank{1cm}\%.
\end{enumerate}

\section*{Round 3: Level 2 Prompts (60\% missing)}
\begin{enumerate}
\item In $\le 4$ lines, state the FSS (1993) mechanism linking \emph{R\&D intensity} to currency hedging. Why is R\&D a sharper proxy for the underinvestment problem than capex?
\item Why should foreign-denominated debt reduce the probability of using FX derivatives? Write the accounting identity implicitly being exploited.
\item Explain why GMS's logit cannot separate ``firm $i$ hedges 1\% of exposure'' from ``firm $i$ hedges 100\% of exposure.'' What estimator would be correct if you had notional amounts (pre- or post-SFAS~119)?
\item Give the logit log-likelihood $\ell(\bm\beta)$ for $n$ independent observations and state the first-order condition.
\end{enumerate}

\section*{Round 4: Minimal Scaffolding (80\% missing)}
From a blank page, reconstruct:
\begin{enumerate}
\item Sample construction rule (what ``ex-ante FX exposure'' operationally means).
\item The full list of proxies with their theoretical mapping.
\item The logit model, marginal effect formula, and reporting convention.
\item The five most important numerical findings (signs and relative magnitudes).
\item Three econometric critiques an AER referee would raise.
\end{enumerate}

\section*{Round 5: Blank Page}
Write, in continuous prose, what GMS (1997) contributes to the risk-management literature, how it differs from Tufano (1996), and why Graham--Rogers (2002) is its natural sequel.

\vspace{6pt}
\begin{drillbox}
\textbf{Speed Drill (60 seconds each).}
\begin{enumerate}
\item Definition of quick ratio and why it proxies for internal liquidity.
\item The sign of R\&D/sales in GMS's logit.
\item Derivation of $\Lambda'(z)=\Lambda(z)(1-\Lambda(z))$.
\item Why ``foreign debt substitutes for FX forwards'' is an identity, not a finding.
\item One reason 41\% usage in 1990 might overstate the economic importance of hedging.
\end{enumerate}
\end{drillbox}

\section*{Quick Reference \& Answer Key}
\begin{answerbox}
\begin{itemize}
\item \textbf{Key formula.} $\Lambda(z)=e^z/(1+e^z)$;\ $\Lambda'(z)=\Lambda(z)(1-\Lambda(z))$; marginal effect at mean $=\bar p(1-\bar p)\beta_k$.
\item \textbf{Quick ratio} $=$ (cash $+$ short-term investments) $/$ current liabilities.
\item \textbf{Headline signs:} R\&D $(+)$, size $(+)$, quick ratio $(-)$, foreign debt $(-)$, leverage $(+)$ but modest, interest coverage $(-)$.
\item \textbf{Mechanism (FSS).} Cash flow correlated with investment opportunities $\Rightarrow$ high-growth-option firms suffer underinvestment in bad states $\Rightarrow$ hedge to keep internal funds available.
\item \textbf{Graham--Smith-style tax convexity.} Not strong in GMS; Graham--Rogers (W07-02) builds the right measure and \emph{rejects} it.
\item \textbf{Classification.} 75--78\%.
\end{itemize}
\end{answerbox}

\begin{keybox}
\textbf{30-second exam pitch.} GMS is the first clean cross-sectional test of why the \emph{typical} large US firm uses FX derivatives. Using 372 Fortune~500 firms with ex-ante FX exposure in 1990, they estimate a logit for the hedge/no-hedge decision. The strongest predictor is R\&D/sales (FSS underinvestment), followed by size (scale economies) and the quick ratio (negative: firms already liquid don't need to hedge). Crucially, foreign-denominated debt is a substitute for FX derivatives, making a pure ``hedging'' reading of derivative use misleading. Logit correctly classifies 75--78\% of firms. Limitations: binary DV ignores intensity; pre-SFAS~119 disclosures are noisy; leverage is treated as exogenous. Graham--Rogers (2002) fixes the last of these by modelling debt and hedging simultaneously.
\end{keybox}

\end{document}
